\section{Limitations and conclusion}
\label{sec:conclusion}
\ours studies instruction-guided executable 3D editing through program, solid, projected-change and conditional response measurements. The completed 3B adaptation improves local strict agreement to 39/252, while its large execution-to-agreement gap identifies remaining edit failures. The flange audit shows how localized visual changes can be missed by whole-shape similarity. Saved geometry and FEM fields make successful illustrative outcomes inspectable.

The model is code-conditioned, and deterministic projections do not establish learned visual perception. The visual audit covers six selected cases, with only four available final 3B outputs; it is not a benchmark-wide metric study. Projected masks miss hidden geometry, depend on rasterization and do not certify feature tolerances. Source-specific target reconstruction, known label defects, local repartitioning and unknown cross-source ancestry further limit the protocol. A single seed and standard QLoRA workflow cannot substantiate a new editing architecture or superiority to published CAD methods.

FEM checks use synthetic contracts on a small subset; they do not establish physical properties, stress convergence or operating safety. Stronger conclusions require matched CAD-editing baselines, repeated training runs, full-test feature/view evaluation and physical tasks with independent analysis contracts. Genuine drawing-conditioned editing remains future work rather than a capability of the current model.
